\documentclass[11pt,a4paper]{article}
\usepackage[margin=2.4cm]{geometry}
\usepackage{amsmath,amssymb,amsthm,bm,booktabs,graphicx,xcolor,mathtools}
\newtheorem{fact}{Fact}
\newtheorem{worked}{Worked example}
\newcommand{\Cof}{\operatorname{cof}}
\newcommand{\tr}{\operatorname{tr}}
\newcommand{\R}{\mathbb{R}}
\newcommand{\SO}{\mathrm{SO}(2)}

\newenvironment{takeaway}
 {\begin{quote}\small\itshape\color{blue!60!black}}
 {\end{quote}}

\title{Polyconvexity, Invariants, and Why We Built Ours This Way\\
\large A didactic companion to the anisotropic PANN construction}
\author{}
\date{\today}

\begin{document}
\maketitle

\begin{abstract}
This is a teaching document, not a paper. Its only goal is for you to
\emph{really} understand, from the ground up, what polyconvexity is,
why hyperelasticity needs this specific and slightly strange-looking
condition instead of ordinary convexity, what the ``three famous
invariants'' $I_1,I_2,I_3$ actually are and why they are polyconvex,
and -- the part that is genuinely specific to this project -- what our
own $15$ features are, why they look the way they do, whether they are
a generalization of $I_1,I_2,I_3$ or something else entirely, and
where the idea came from. Every numerical example below was computed
and checked, not invented, so you can reproduce every number yourself.
\end{abstract}

\tableofcontents

\section{Why do we care about convexity at all?}
\label{sec:motivation}

Start from the most basic possible question, with no elasticity in it
yet: why is convexity such a big deal in optimization and in physics?

\begin{takeaway}
A convex function has no ``fake'' local minima to get stuck in, and a
line connecting any two points on its graph always stays above the
graph. This one geometric fact is what makes minimization problems
built from convex functions well-behaved: gradient-based solvers
(including Newton's method) converge reliably, and a minimizer is
guaranteed to exist under mild extra conditions.
\end{takeaway}

Take the simplest possible example, $f(x)=x^2$. It has one minimum, at
$x=0$, and Newton's method finds it in one step from anywhere. Now
compare $g(x)=x^4-3x^2$ (a double-well shape): it has \emph{two} local
minima and a local maximum in between. A gradient-based solver started
near the ``wrong'' well will happily converge to a minimizer that is
not the global one, and if you are minimizing an energy that is
supposed to represent physical equilibrium, converging to the wrong
local minimum is a real, physical wrong answer, not just a numerical
inconvenience.

In hyperelasticity, the ``optimization problem'' is: find the
displacement field that minimizes the total elastic energy, subject to
the boundary conditions. If the energy density $W$ (as a function of
the deformation) is convex, classical results (the \emph{direct method
of the calculus of variations}) give you existence of a minimizer more
or less for free, and Newton's method on the discretized problem tends
to behave itself. This is the entire motivation for wanting some kind
of convexity. The question is: convexity of \emph{what}, exactly?

\section{The naive answer fails: $W$ cannot be convex in $F$}
\label{sec:naive-fails}

The most naive thing to ask for is that $W$, as a function of the
deformation gradient $F$ (a $2\times2$ matrix, in our 2D setting), be
convex in $F$. This fails, and it fails for a clean, checkable reason
that has nothing to do with any specific material: \emph{objectivity}
(frame indifference) makes it impossible.

Objectivity says $W(QF)=W(F)$ for every rotation $Q\in\SO$: rotating
the deformed body rigidly, without straining it further, cannot change
the stored energy. Take $Q=-I$ (a $180^\circ$ rotation, which is a
genuine element of $\SO$ in two dimensions). Objectivity forces
\[
 W(-F)=W(F).
\]
Now if $W$ were convex in $F$, the value at the midpoint of the segment
from $F$ to $-F$ would have to be no larger than the average of the
endpoints:
\[
 W\!\left(\tfrac12F+\tfrac12(-F)\right) = W(\bm 0) \;\leq\; \tfrac12W(F)+\tfrac12W(-F) = W(F).
\]
But the midpoint is the \emph{zero matrix}, $F=\bm 0$, which has
$\det F=0$: physically, this is a state where the material has been
crushed to zero volume, and any physically sensible energy must blow
up there ($W\to\infty$), not stay finite and below $W(F)$. Convexity in
$F$ and objectivity are flatly incompatible. This is not a defect of
any particular model; it is a theorem (it goes back to Coleman and
Noll, and is discussed in essentially every modern treatment,
e.g.\ Ciarlet's textbook~\cite{Ciarlet1988}).

\begin{takeaway}
You cannot require $W$ to be an ordinary convex function of $F$. Any
usable notion of convexity for elasticity has to be something weaker,
compatible with $W(-F)=W(F)$ and with $W\to\infty$ as the material is
crushed.
\end{takeaway}

\section{Ball's trick: lift the problem into a bigger space}
\label{sec:lifting-trick}

Here is the central idea, and it is worth understanding it once, in
the abstract, completely divorced from elasticity, because then the
elasticity version stops looking mysterious.

\begin{takeaway}
A function can fail to be convex in its \emph{original} variable, and
yet become convex the moment you feed it some \emph{extra, derived}
quantity as if that quantity were an independent input.
\end{takeaway}

\begin{worked}[The lifting trick, with no elasticity in it]
\label{ex:lifting}
Let $F=\begin{psmallmatrix}a&b\\c&d\end{psmallmatrix}$ and consider the
function $h(F)=\det F=ad-bc$. As a function of the four numbers
$(a,b,c,d)$, $h$ is a saddle-shaped bilinear form: its Hessian is
indefinite (it has both positive and negative eigenvalues), so $h$ is
\emph{not} convex in $F$. Nudging $b$ and $c$ in the right directions
can decrease $h$ below the average of two points where it is large --
exactly what convexity forbids.

Now play a trick: introduce a new, extra symbol $\delta$, declare
$\delta:=\det F$, and consider the function of the \emph{pair}
$(F,\delta)$ given by
\[
 \tilde h(F,\delta) := \delta.
\]
This is now trivially convex (it is even affine/linear) in the pair
$(F,\delta)$, because it does not depend on $F$ at all once $\delta$ is
treated as a free-standing coordinate! Of course $\tilde h(F,\det F) =
h(F)$ recovers the original, non-convex function -- nothing physical
changed. What changed is that we are now allowed to ask a
\emph{different, easier} question: is $\tilde h$ jointly convex in the
enlarged pair of independent variables $(F,\delta)$, rather than in $F$
alone along the curved, constrained set $\{(F,\det F)\}$? The answer
can be yes even when the original question's answer is no.
\end{worked}

This is, almost verbatim, what Ball~\cite{Ball1977} does for
elasticity. Instead of asking whether $W$ is convex in $F$, he asks
whether there exists a function $G$ of \emph{three independent block
variables} $(F,\Cof F,\det F)$ -- treating the cofactor matrix $\Cof F$
and the determinant $\det F$ as if they were extra, free-standing
inputs, exactly like $\delta$ above -- such that
\begin{equation}
 W(F)=G(F,\Cof F,\det F),\qquad G\text{ jointly convex on }
 \R^{2\times2}\times\R^{2\times2}\times\R.
 \label{eq:polyconvex-def}
\end{equation}
This is \emph{polyconvexity}. It is strictly weaker than convexity in
$F$ (every $F$-convex function is trivially polyconvex, by ignoring the
extra blocks, but not conversely), it is compatible with objectivity
(because $\Cof(-F)=\Cof F$ and $\det(-F)=\det F$ in 2D, so
$G(-F,\Cof F,\det F)$ need not equal $G(F,\Cof F,\det F)$ for the same
reason $\delta$ didn't care about the sign flip trick above -- the
non-convex, sign-sensitive part of the problem is exactly the part that
got ``absorbed'' into the extra blocks), and it is exactly the
condition Ball needed, together with a growth condition, to prove
existence of energy minimizers by the direct method. It says nothing,
by itself, about the sign of $\partial^2W/\partial F^2$ at a given
deformation -- that is a completely different, and strictly stronger,
question, discussed in the main paper.

\section{The three famous invariants}
\label{sec:classical-invariants}

Now to your direct question: yes, you have heard of three famous
invariants, and here they are. For the right Cauchy--Green tensor
$C=F^TF$ (symmetric, positive definite), the classical invariants used
in essentially every hyperelasticity textbook (see
Ogden~\cite{Ogden1984}, or Ciarlet~\cite{Ciarlet1988}) are
\begin{align}
 I_1 &= \tr C, \label{eq:I1}\\
 I_2 &= \tfrac12\big[(\tr C)^2-\tr(C^2)\big]
       = \tr\!\big(\Cof(C)\big)\quad\text{(in 3D)}, \label{eq:I2}\\
 I_3 &= \det C = J^2. \label{eq:I3}
\end{align}
Physically: $I_1$ is (twice) the sum of the squared principal
stretches, so it measures how much the material has been stretched
overall, in all directions combined; $I_3=J^2$ measures the squared
change in volume; $I_2$ is related to how much \emph{areas} (rather
than lengths) have changed, which is why it is built from $\Cof(C)$,
the object that maps how oriented area elements transform under $F$.

The two most famous compressible/incompressible hyperelastic models in
existence are built directly from these:
\[
 W_{\rm neo\text{-}Hookean}=c_1(I_1-3), \qquad
 W_{\rm Mooney\text{-}Rivlin}=c_1(I_1-3)+c_2(I_2-3),
\]
and both are polyconvex, for a clean structural reason you can already
see: $I_1=\tr C=|F|^2$ (a sum of squares of the entries of $F$) is
convex in $F$ directly -- it lives entirely in the \emph{first} block
of~\eqref{eq:polyconvex-def} -- and $I_2$, being built from $\Cof(C)$,
lives in the \emph{second} block. For non-negative coefficients $c_1,
c_2\geq0$, a non-negative combination of functions that are each convex
in their own block is again jointly convex in the combined variable
$(F,\Cof F,J)$ -- exactly the ``sum of convex pieces is convex'' rule
used throughout the main paper.

\begin{fact}[A 2D subtlety worth knowing]
\label{fact:2d-degeneracy}
In three dimensions, $I_1,I_2,I_3$ are three genuinely independent
numbers. In two dimensions -- our setting -- $C$ has only two
eigenvalues $\lambda_1,\lambda_2$ (squared principal stretches), and
the ``$I_2$'' role (sum of products of pairs of eigenvalues) collapses
onto a single pair, $\lambda_1\lambda_2=\det C=I_3$. So in 2D, the
classical $I_2$ and $I_3$ are literally the same number, and the naive
``three invariants'' picture only has two independent scalars,
$I_1$ and $J$. This is exactly why Ball's 2D theory is phrased using
the whole \emph{matrix} $\Cof F$ as a block, not just its scalar trace:
treating $\Cof F$ as a full $2\times2$ object (not merely
``the second invariant'') leaves room to extract genuinely new,
non-isotropic information from it, which a single scalar trace could
never do (its trace, as just noted, only gives you $I_1$ back). This is
precisely the room our own construction below uses.
\end{fact}

\section{Are we using $I_1,I_2,I_3$? -- Yes, generalized}
\label{sec:are-we-using-them}

Here is the direct answer to ``are we using the three famous
invariants,'' made precise rather than hand-wavy.

\begin{fact}[$I_1$ is a very special case of our own construction]
\label{fact:i1-special-case}
A basic, exact identity: for \emph{any} orthonormal basis
$\{e_1,e_2\}$ of $\R^2$ (not just the standard one),
\[
 I_1=\tr C = e_1^TCe_1+e_2^TCe_2 .
\]
That is, $I_1$ is \emph{already}, secretly, a weighted sum of
$d\cdot Cd$ over two chosen unit directions, with weight $1$ on each
and power $p=1$. Compare this to our own feature family (main paper,
Eq.~for $Q^F_{k,p}$),
\[
 Q_{k,p}^F=\sum_{i=1}^3 w_{ki}\,(d_{ki}\cdot Cd_{ki})^p .
\]
$I_1$ is exactly the special case of this formula with: only
\emph{two} directions instead of three, those two directions
\emph{orthogonal} to each other, weights \emph{equal to} $1$, and power
$p=1$. Our construction is not analogous to $I_1$; it is a literal,
direct generalization of it -- more directions (three, not two), no
requirement that they be orthogonal, general positive weights (subject
to the balance condition below, which is exactly what replaces
``orthonormal'' once you give up orthogonality), and higher powers
$p=2,3$ instead of $p=1$.
\end{fact}

The same story holds one level up: our $Q_{k,p}^H$ features use
$d\cdot\Cof(C)d$ in exactly the same way $Q^F$ uses $d\cdot Cd$, so
they generalize the $\Cof$-based ($I_2$-flavoured) invariant in exactly
the same sense. And $J,J^2$ are not generalized at all -- they are
$I_3$ and $I_3^2$, used completely unchanged, because the volumetric
invariant does not need an anisotropic generalization: volume change is
already a single, direction-independent scalar, and there is nothing
anisotropic to generalize.

\begin{takeaway}
So: yes, we are using $I_1$, $I_2$ (via $\Cof$), and $I_3=J^2$ -- but
not as three bare numbers. We use \emph{many} directional versions of
the $I_1$-type and $I_2$-type constructions (three direction systems,
each contributing terms at powers $p=2,3$), because a material with
\emph{no assumed symmetry} needs to be able to respond differently in
different directions, and a single scalar $I_1$ or $I_2$, being built
from an orthonormal (or in fact \emph{any}) basis by summation, is
completely blind to direction by construction.
\end{takeaway}

\section{Why $I_1,I_2,I_3$ alone cannot describe an anisotropic material}
\label{sec:isotropy-blindness}

This needs to be seen numerically to really land, so here it is,
computed for real, not asserted.

\begin{worked}[Two different materials that $I_1$ and $J$ cannot tell apart]
\label{ex:isotropy-blind}
Take two deformation states,
\[
 F_a=\begin{psmallmatrix}1.4&0\\0&0.9\end{psmallmatrix}
 \quad(\text{stretch }1.4\times\text{ along }x,\ \text{compress to }0.9\times\text{ along }y),
 \qquad
 F_b=\begin{psmallmatrix}0.9&0\\0&1.4\end{psmallmatrix}
 \quad(\text{roles of }x,y\text{ swapped}).
\]
These are physically very different states for an anisotropic
material -- one stretches hard along $x$ and compresses along $y$, the
other does the opposite -- and a material whose stiffness depends on
direction should, in general, store a different amount of energy in
each. Yet:
\[
 I_1(F_a)=I_1(F_b)=2.77, \qquad J(F_a)=J(F_b)=1.26 .
\]
Exactly equal, to every decimal place, because trace and determinant
literally do not know which axis is which. \emph{Any} energy built
purely from $I_1,I_2,I_3$ (i.e., any \emph{isotropic} hyperelastic
model, no matter how the invariants are combined) is mathematically
incapable of distinguishing $F_a$ from $F_b$. This is not a
approximation error; it is exact and structural.
\end{worked}

\section{Our construction: direction systems, balance, and why the cubic power is not decorative}
\label{sec:our-construction}

To fix Example~\ref{ex:isotropy-blind}, you need at least one feature
that is \emph{not} blind to direction. That is what the three
direction systems (Table of direction systems in the main paper) are
for. Each system $k$ consists of three unit vectors $d_{k1},d_{k2},
d_{k3}$ at specific angles, with specific positive weights $w_{k1},
w_{k2},w_{k3}$, satisfying the \emph{balance condition}
\begin{equation}
 \sum_{i=1}^3 w_{ki}\,d_{ki}\otimes d_{ki} = I .
 \label{eq:balance}
\end{equation}
This condition is the direct 2D generalization of ``orthonormal with
unit weights'': it says the chosen (possibly oblique, non-orthogonal)
directions, weighted appropriately, still add up to the identity, the
same way $e_1\otimes e_1+e_2\otimes e_2=I$ does for an orthonormal
pair. This is exactly what makes the reference-state stress-free
proof (C4 in the main paper) go through with a single scalar
correction, and it is why the model behaves \emph{isotropically at the
undeformed state} even though it is fully anisotropic away from it --
a real material with a texture or weave usually looks locally isotropic
right at rest and only reveals its anisotropy once you deform it, which
is exactly the behavior this construction reproduces by design.

\begin{worked}[Computing $Q^F_{1,2}$ and $Q^F_{1,3}$ by hand, for real]
\label{ex:worked-full}
System~1 uses angles $0^\circ,50^\circ,130^\circ$ with weights
$0.2959,0.8520,0.8520$. For $F_a$ above, $C=\mathrm{diag}(1.96,0.81)$.
The three directions are $d_1=(1,0)$, $d_2=(0.6428,0.7660)$,
$d_3=(-0.6428,0.7660)$, giving
\[
 d_1\cdot Cd_1=1.96,\qquad d_2\cdot Cd_2=d_3\cdot Cd_3=1.2852 .
\]
So
\[
 Q^F_{1,2}(F_a)=0.2959(1.96)^2+2\times0.8520(1.2852)^2=3.9511,
\]
\[
 Q^F_{1,3}(F_a)=0.2959(1.96)^3+2\times0.8520(1.2852)^3=5.8449.
\]
Repeating for $F_b$ (swap the roles of $x,y$, so effectively
$d\cdot Cd$ picks up the $0.81$ and $1.96$ in the opposite pattern)
gives
\[
 Q^F_{1,2}(F_b)=3.95107,\qquad Q^F_{1,3}(F_b)=5.7357.
\]
Now compare the two states directly:
\[
 \big|Q^F_{1,2}(F_a)-Q^F_{1,2}(F_b)\big| \approx 0.00001,
 \qquad
 \big|Q^F_{1,3}(F_a)-Q^F_{1,3}(F_b)\big| \approx 0.109 .
\]
This is the single most important number in this whole document.
\textbf{The power-2 feature barely tells $F_a$ and $F_b$ apart}
(a difference in the fifth decimal, for this system) -- it turns out a
balanced \emph{quadratic} aggregate in 2D obeys an almost-exact hidden
symmetry under swapping the two axes, so it is a nearly-blind, weak
anisotropy detector, not so different from $I_1$ itself in this
respect. \textbf{The power-3 feature, by contrast, differs by about
$1.9\%$ of its own value} -- a real, structural, non-accidental
difference. This is exactly why the model uses both $p=2$ and $p=3$
terms for every direction system, and why the main paper states this
is ``not cosmetic'': without the cubic terms, the model would be far
closer to isotropic than intended, almost by accident, for reasons that
have nothing to do with training and everything to do with a 2D
geometric coincidence at even powers.
\end{worked}

Why \emph{three} separate direction systems, rather than one? One
system of three directions already breaks isotropy (Example
\ref{ex:worked-full} proves that), but it only samples the material's
directional response along three specific angles. A real anisotropic
material's stiffness can vary continuously with direction in a
complicated way; three independent, differently-angled systems (nine
directions in total, each individually balanced) give the model many
more independent ``probes'' of directional stiffness to combine,
exactly as using more sensors at more angles gives a richer picture of
an unknown directional signal than a single triple of sensors would.

\section{Where did this come from -- did we invent it?}
\label{sec:provenance}

Split honestly into two layers, because the honest answer has two
different parts.

\paragraph{The overall strategy (not ours).} Building a polyconvex
energy as an input-convex neural network (ICNN) applied to hand-picked,
individually-polyconvex invariants is Klein et al.'s
approach~\cite{Klein2022}, itself sitting on top of Ball's original
existence theory~\cite{Ball1977} and the standard invariant-based
hyperelasticity tradition going back to Ogden and
earlier~\cite{Ogden1984}. Klein et al.\ tie their invariants to a
\emph{named} material symmetry group (cubic, transversely isotropic,
\dots) via a classical structural tensor -- a well-established, textbook
technique (structural-tensor / integrity-basis representation theory
for anisotropic invariants) applied by many authors before them.

\paragraph{The specific mechanism (ours).} What is genuinely this
project's own construction is the decision \emph{not} to tie the
invariants to any named symmetry group at all: instead, use several
generic direction systems, individually balanced
via~\eqref{eq:balance}, with no $90^\circ$-rotation closure and no
$D_4$ averaging built in anywhere. This is the specific answer to the
open problem Linden et al.~\cite{Linden2023} flag explicitly in their
own paper -- that a complete, polyconvexity-compatible invariant set
may not exist for a general, unknown symmetry class -- and it is what
lets this construction claim exact objectivity, an exact stress-free
reference state, and exact polyconvexity \emph{without ever assuming or
checking} which symmetry class, if any, the material belongs to. The
specific numerical angles and weights in the direction-systems table
are one valid solution of the (underdetermined) balance
condition~\eqref{eq:balance} -- there are infinitely many valid
choices, and the ones used were not derived from any deeper physical
principle beyond satisfying~\eqref{eq:balance} to machine precision and
providing three visibly different, non-$D_4$-related angular
patterns; if you want a cleaner or differently-optimized choice, the
balance condition is the only real constraint you must preserve.

\section{A short recap, in one paragraph}
\label{sec:recap}

Ordinary convexity in $F$ is impossible for any objective energy
(Section~\ref{sec:naive-fails}). Ball's fix, polyconvexity, asks for
convexity not in $F$ but in the enlarged, ``lifted'' triple
$(F,\Cof F,\det F)$, treated as independent variables
(Section~\ref{sec:lifting-trick}) -- exactly the same trick as
promoting $\det F$ to its own free coordinate in
Example~\ref{ex:lifting}. The classical invariants $I_1,I_2,I_3$ are
the textbook way to build polyconvex energies out of this triple
(Section~\ref{sec:classical-invariants}), but they are isotropic by
construction and provably blind to direction
(Example~\ref{ex:isotropy-blind}). Our own features are a direct,
literal generalization of $I_1$ and the $\Cof$-based invariant --
more directions, general weights satisfying a balance condition, and
crucially both quadratic \emph{and cubic} powers, because the quadratic
power alone is a nearly-blind anisotropy detector in 2D
(Example~\ref{ex:worked-full}). The overall ICNN-on-invariants strategy
is Klein et al.'s; the specific choice to use generic, symmetry-free
direction systems instead of a named structural tensor is this
project's own contribution, built to answer an open problem stated
explicitly in the literature it responds to.

\begin{thebibliography}{9}
\bibitem{Ball1977}
J.~M. Ball. Convexity conditions and existence theorems in nonlinear
elasticity. \emph{Archive for Rational Mechanics and Analysis}, 63:337--403,
1977.
\bibitem{Ciarlet1988}
P.~G. Ciarlet. \emph{Mathematical Elasticity, Volume I: Three-Dimensional
Elasticity}. North-Holland, 1988.
\bibitem{Ogden1984}
R.~W. Ogden. \emph{Non-Linear Elastic Deformations}. Ellis Horwood,
Chichester, 1984 (reprinted by Dover, 1997).
\bibitem{Klein2022}
D.~K. Klein, M. Fern\'andez, R.~J. Martin, P. Neff, and O. Weeger.
Polyconvex anisotropic hyperelasticity with neural networks.
\emph{Journal of the Mechanics and Physics of Solids}, 159:104703, 2022.
\bibitem{Linden2023}
L. Linden, D.~K. Klein, K.~A. Kalina, J. Brummund, O. Weeger, and
M. K\"astner. Neural networks meet hyperelasticity: a guide to enforcing
physics. \emph{Journal of the Mechanics and Physics of Solids},
179:105363, 2023.
\end{thebibliography}

\end{document}
